% problem_id: S001-lemma-cohomology-on-subcategory
% source_id: S001 (Stacks Project)
% source_file: stacks-sheaves.tex
% statement_locator: stacks-sheaves.tex lines 4218--4243
% proof_locator: stacks-sheaves.tex lines 4245--4401
% env: lemma
% id_note: label 在本来源内唯一
% pairing: tight (verified)
% status: complete (statement + author proof verbatim + reason + locators)
%
\begin{lemma}
\label{lemma-cohomology-on-subcategory}
Let $S$ be a scheme. Let $\mathcal{X}$ be an algebraic stack over $S$.
Let $\tau = \etale$ (resp.\ $\tau = fppf$). Let
$\mathcal{X}' \subset \mathcal{X}$ be a full subcategory with the
following properties
\begin{enumerate}
\item if $x \to x'$ is a morphism of $\mathcal{X}$ which lies over a
smooth (resp.\ flat and locally finitely presented) morphism of
schemes and $x' \in \Ob(\mathcal{X}')$, then $x \in \Ob(\mathcal{X}')$, and
\item there exists an object $x \in \Ob(\mathcal{X}')$ lying over
a scheme $U$ such that the associated $1$-morphism
$x : (\Sch/U)_{fppf} \to \mathcal{X}$ is smooth and surjective.
\end{enumerate}
We get a site $\mathcal{X}'_\tau$ by declaring a covering of $\mathcal{X}'$
to be any family of morphisms $\{x_i \to x\}$ in $\mathcal{X}'$ which is a
covering in $\mathcal{X}_\tau$. Then the inclusion functor
$\mathcal{X}' \to \mathcal{X}_\tau$ is fully faithful, cocontinuous, and
continuous, whence defines a morphism of topoi
$$
g : \Sh(\mathcal{X}'_\tau) \longrightarrow \Sh(\mathcal{X}_\tau)
$$
and $H^p(\mathcal{X}'_\tau, g^{-1}\mathcal{F}) =
H^p(\mathcal{X}_\tau, \mathcal{F})$ for all $p \geq 0$ and all
$\mathcal{F} \in \textit{Ab}(\mathcal{X}_\tau)$.
\end{lemma}

\begin{proof}
Note that assumption (1) implies that if $\{x_i \to x\}$ is a covering
of $\mathcal{X}_\tau$ and $x \in \Ob(\mathcal{X}')$, then we have
$x_i \in \Ob(\mathcal{X}')$. Hence we see that $\mathcal{X}' \to \mathcal{X}$
is continuous and cocontinuous as the coverings of objects of
$\mathcal{X}'_\tau$ agree with their coverings seen as objects of
$\mathcal{X}_\tau$. We obtain the morphism $g$ and the functor
$g^{-1}$ is identified with the restriction functor, see
Sites, Lemma \ref{sites-lemma-when-shriek}.

\medskip\noindent
In particular, if $\{x_i \to x\}$ is a covering in $\mathcal{X}'_\tau$,
then for any abelian sheaf $\mathcal{F}$ on $\mathcal{X}$ then
$$
\check H^p(\{x_i \to x\}, g^{-1}\mathcal{F}) =
\check H^p(\{x_i \to x\}, \mathcal{F})
$$
Thus if $\mathcal{I}$ is an injective abelian sheaf on $\mathcal{X}_\tau$
then we see that the higher {\v C}ech cohomology groups are zero
(Cohomology on Sites,
Lemma \ref{sites-cohomology-lemma-injective-trivial-cech}).
Hence $H^p(x, g^{-1}\mathcal{I}) = 0$ for all objects $x$
of $\mathcal{X}'$
(Cohomology on Sites,
Lemma \ref{sites-cohomology-lemma-cech-vanish-collection}).
In other words injective abelian sheaves on $\mathcal{X}_\tau$
are right acyclic for the functor $H^0(x, g^{-1}-)$.
It follows that $H^p(x, g^{-1}\mathcal{F}) = H^p(x, \mathcal{F})$
for all $\mathcal{F} \in \textit{Ab}(\mathcal{X})$ and all
$x \in \Ob(\mathcal{X}')$.

\medskip\noindent
Choose an object $x \in \mathcal{X}'$ lying over a scheme $U$
as in assumption (2). In particular $\mathcal{X}/x \to \mathcal{X}$
is a morphism of algebraic stacks which representable by algebraic spaces,
surjective, and smooth. (Note that $\mathcal{X}/x$ is equivalent to
$(\Sch/U)_{fppf}$, see Lemma \ref{lemma-localizing}.)
The map of sheaves
$$
h_x \longrightarrow *
$$
in $\Sh(\mathcal{X}_\tau)$ is surjective. Namely, for any object $x'$
of $\mathcal{X}$ there exists a $\tau$-covering $\{x'_i \to x'\}$
such that there exist morphisms $x'_i \to x$, see
Lemma \ref{lemma-surjective-flat-locally-finite-presentation}.
Since $g$ is exact, the map of sheaves
$$
g^{-1}h_x \longrightarrow * = g^{-1}*
$$
in $\Sh(\mathcal{X}'_\tau)$ is surjective also. Let $h_{x, n}$ be
the $(n + 1)$-fold product $h_x \times \ldots \times h_x$.
Then we have spectral sequences
\begin{equation}
\label{equation-spectral-sequence-one}
E_1^{p, q} = H^q(h_{x, p}, \mathcal{F}) \Rightarrow
H^{p + q}(\mathcal{X}_\tau, \mathcal{F})
\end{equation}
and
\begin{equation}
\label{equation-spectral-sequence-two}
E_1^{p, q} = H^q(g^{-1}h_{x, p}, g^{-1}\mathcal{F}) \Rightarrow
H^{p + q}(\mathcal{X}'_\tau, g^{-1}\mathcal{F})
\end{equation}
see Cohomology on Sites,
Lemma \ref{sites-cohomology-lemma-cech-to-cohomology-sheaf-sets}.

\medskip\noindent
Case I: $\mathcal{X}$ has a final object $x$ which is also an object of
$\mathcal{X}'$. This case follows immediately from the discussion
in the second paragraph above.

\medskip\noindent
Case II: $\mathcal{X}$ is representable by an algebraic space $F$.
In this case the sheaves $h_{x, n}$ are representable by an
object $x_n$ in $\mathcal{X}$. (Namely, if $\mathcal{S}_F = \mathcal{X}$
and $x : U \to F$ is the given object, then $h_{x, n}$ is representable
by the object $U \times_F \ldots \times_F U \to F$ of $\mathcal{S}_F$.)
It follows that $H^q(h_{x, p}, \mathcal{F}) = H^q(x_p, \mathcal{F})$.
The morphisms $x_n \to x$ lie over smooth morphisms of schemes, hence
$x_n \in \mathcal{X}'$ for all $n$. Hence
$H^q(g^{-1}h_{x, p}, g^{-1}\mathcal{F}) = H^q(x_p, g^{-1}\mathcal{F})$.
Thus in the two spectral sequences
(\ref{equation-spectral-sequence-one}) and
(\ref{equation-spectral-sequence-two}) above the $E_1^{p, q}$ terms agree
by the discussion in the second paragraph. The lemma follows in Case II
as well.

\medskip\noindent
Case III: $\mathcal{X}$ is an algebraic stack. We claim that in this case
the cohomology groups $H^q(h_{x, p}, \mathcal{F})$ and
$H^q(g^{-1}h_{x, n}, g^{-1}\mathcal{F})$ agree by Case II above.
Once we have proved this the result will follow as before.

\medskip\noindent
Namely, consider the category $\mathcal{X}/h_{x, n}$, see
Sites, Lemma \ref{sites-lemma-localize-topos-site}.
Since $h_{x, n}$ is the $(n + 1)$-fold product of $h_x$ an
object of this category is an $(n + 2)$-tuple
$(y, s_0, \ldots, s_n)$ where $y$ is an object of $\mathcal{X}$ and each
$s_i : y \to x$ is a morphism of $\mathcal{X}$.
This is a category over $(\Sch/S)_{fppf}$. There is an equivalence
$$
\mathcal{X}/h_{x, n}
\longrightarrow
(\Sch/U)_{fppf} \times_\mathcal{X} \ldots \times_\mathcal{X} (\Sch/U)_{fppf}
=: \mathcal{U}_n
$$
over $(\Sch/S)_{fppf}$. Namely, if $x : (\Sch/U)_{fppf} \to \mathcal{X}$ also
denotes the $1$-morphism associated with $x$ and
$p : \mathcal{X} \to (\Sch/S)_{fppf}$ the structure functor,
then we can think of $(y, s_0, \ldots, s_n)$ as
$(y, f_0, \ldots, f_n, \alpha_0, \ldots, \alpha_n)$
where $y$ is an object of $\mathcal{X}$, $f_i : p(y) \to p(x)$ is a
morphism of schemes, and $\alpha_i : y \to x(f_i)$ an isomorphism.
The category of $2n+3$-tuples
$(y, f_0, \ldots, f_n, \alpha_0, \ldots, \alpha_n)$
is an incarnation of the $(n + 1)$-fold fibred product $\mathcal{U}_n$
of algebraic stacks displayed above, as we discussed in
Section \ref{section-cech}.
By Cohomology on Sites, Lemma
\ref{sites-cohomology-lemma-cohomology-on-sheaf-sets}
we have
$$
H^p(\mathcal{U}_n, \mathcal{F}|_{\mathcal{U}_n}) =
H^p(\mathcal{X}/h_{x, n}, \mathcal{F}|_{\mathcal{X}/h_{x, n}}) =
H^p(h_{x, n}, \mathcal{F}).
$$
Finally, we discuss the ``primed'' analogue of this. Namely,
$\mathcal{X}'/h_{x, n}$ corresponds, via the equivalence above
to the full subcategory $\mathcal{U}'_n \subset \mathcal{U}_n$
consisting of those tuples
$(y, f_0, \ldots, f_n, \alpha_0, \ldots, \alpha_n)$
with $y \in \mathcal{X}'$. Hence certainly property (1) of the
statement of the lemma holds
for the inclusion $\mathcal{U}'_n \subset \mathcal{U}_n$.
To see property (2) choose an object $\xi = (y, s_0, \ldots, s_n)$ which
lies over a scheme $W$ such that $(\Sch/W)_{fppf} \to \mathcal{U}_n$
is smooth and surjective (this is possible as $\mathcal{U}_n$ is
an algebraic stack). Then
$(\Sch/W)_{fppf} \to \mathcal{U}_n \to (\Sch/U)_{fppf}$
is smooth as a composition of base changes of the morphism
$x : (\Sch/U)_{fppf} \to \mathcal{X}$, see
Algebraic Stacks, Lemmas
\ref{algebraic-lemma-base-change-representable-transformations-property} and
\ref{algebraic-lemma-composition-representable-transformations-property}.
Thus axiom (1) for $\mathcal{X}$ implies that $y$ is an object of
$\mathcal{X}'$ whence $\xi$ is an object of $\mathcal{U}'_n$.
Using again
$$
H^p(\mathcal{U}'_n, \mathcal{F}|_{\mathcal{U}'_n}) =
H^p(\mathcal{X}'/h_{x, n}, \mathcal{F}|_{\mathcal{X}'/h_{x, n}}) =
H^p(g^{-1}h_{x, n}, g^{-1}\mathcal{F}).
$$
we now can use Case II for 
$\mathcal{U}'_n \subset \mathcal{U}_n$
to conclude.
\end{proof}
